% ====================================================================
% Section: Convergence of independent classical-physics cores.
% The new centerpiece. Faithful to data/equivalence_{newton,archimedes,
% pascal}.json, all measured on the identical CONGA-K16 / DRILL / TOST bench.
% ====================================================================
\section{Convergence: Three Classical Physics, One Destination}
\label{sec:convergence}

If a single physically-motivated rule matched engineered routing, it could be a
coincidence---one lucky analogy. The stronger claim, and the one this section
supports, is that \emph{several independent classical-physics principles, each
formalised on its own and sharing no mechanism, converge on the same operating
point as engineered load balancers}. Convergence from independent starting
points is evidence that the result reflects the structure of the problem, not
the cleverness of one analogy.

\subsection{Three cores, three mechanisms}

We implement three routing cores. All share one search engine---an exact
hop-budget dynamic program descending a per-edge potential---and differ
\emph{only} in the added physical term. This makes the comparison fair by
construction: any difference between cores comes from the physics, not the
search.

\begin{itemize}
\item \textbf{Newton (reaction).} A dropped packet deposits a decaying scar on
the offending edge; later packets feel a repulsive reaction and route away. The
mechanism is \emph{reactive} and \emph{temporal}: it responds to losses already
incurred, and the scar fades when not refreshed. (This is the two-term core of
Section~\ref{sec:core}.)

\item \textbf{Archimedes (anticipation).} Each edge exerts a buoyant push that
grows with the square of its congestion density above a flotation threshold
$\rho_0$: $g\cdot\max(0,\rho-\rho_0)^2$. The mechanism is \emph{anticipatory}
and \emph{non-linear}: a packet is repelled from an edge \emph{before} it drops
a packet, and only once density crosses the threshold. It carries no memory and
no packet state.

\item \textbf{Pascal (diffusion).} An edge inherits a share of the congestion of
its graph neighbours, so a packet avoids not just congested edges but congested
\emph{regions}. The mechanism is \emph{spatial}: pressure transmits to adjacent
edges, as in a confined fluid.
\end{itemize}

These are not variants of one idea. Reaction-to-the-past, anticipation-of-the-
present, and spatial-diffusion are physically distinct, and share no term beyond
the common distance-plus-congestion substrate that any router needs.

\subsection{All three reach parity}

We evaluate each core on the identical bench used throughout: fifteen
topologies, twenty seeds, flow-level simulation, statistical equivalence by TOST
at margin $\delta = 2$ percentage points, against DRILL and against CONGA at a
\emph{fair} path budget ($K=16$). Table~\ref{tab:convergence} reports the
equivalence counts and, more tellingly, the per-topology pattern.

\begin{table}[t]
\centering
\caption{Equivalence (TOST, $\delta=2$pp) of three independent physical cores
with engineered routers, on the identical bench. ``C'' marks equivalence with
CONGA$_{K=16}$, ``D'' with DRILL. The cores agree cell by cell far more than
they differ, and fail on the same topologies.}
\label{tab:convergence}
\small
\begin{tabular}{lccc}
\toprule
Topology & Newton & Archimedes & Pascal \\
\midrule
Grid $5$     & \phantom{C}D & CD           & \phantom{C}D \\
Grid $6$     & \phantom{C}D & CD           & \phantom{C}D \\
Grid $7$     & \phantom{C}D & CD           & CD \\
Grid $8$     & CD           & CD           & CD \\
Grid $10$    & CD           & C\phantom{D} & CD \\
Grid $12$    & CD           & CD           & CD \\
WS $n30\,k4$ & --           & --           & -- \\
WS $n50\,k4$ & C\phantom{D} & C\phantom{D} & C\phantom{D} \\
WS $n50\,k6$ & CD           & CD           & CD \\
WS $n80\,k4$ & C\phantom{D} & C\phantom{D} & C\phantom{D} \\
BA $n50\,m2$ & CD           & CD           & CD \\
BA $n50\,m3$ & CD           & CD           & CD \\
BA $n80\,m2$ & CD           & CD           & CD \\
G\'EANT      & --           & \phantom{C}D & -- \\
Abilene      & --           & \phantom{C}D & -- \\
\midrule
\textbf{$\equiv$ CONGA$_{K16}$} & 9/15 & 12/15 & 10/15 \\
\textbf{$\equiv$ DRILL}         & 10/15 & 11/15 & 10/15 \\
\bottomrule
\end{tabular}
\end{table}

Two features of the table matter more than the totals. First, the three cores
land in the same band---9 to 12 of 15 against each baseline---despite their
mechanisms having nothing in common. Second, and more striking, \emph{they agree
cell by cell}: where one core reaches equivalence the others tend to, and where
one fails the others fail. The disagreements are few and small.

\subsection{They fail in the same places}

The shared failures are the most informative part of the result. All three cores
miss on \textbf{WS\,$n30\,k4$}, a dense small-world graph where DRILL's per-hop
reactivity is genuinely hard to match, and all three miss (against CONGA) on the
two real WAN backbones, \textbf{G\'EANT} and \textbf{Abilene}. These backbones
have low tube/sp ratio---little room to manoeuvre among near-shortest
paths---and Section~\ref{sec:applicability} shows the same metric predicts the
failures of all three cores, not just one.

That independent physical mechanisms succeed on the same topologies and fail on
the same topologies is the core finding. It means the boundary of the result is
a property of the \emph{problem}---the topology's room to manoeuvre---rather than
of any single physical analogy. Load distribution on a graph with sufficient
near-shortest diversity is recoverable by classical physics; which physics one
chooses barely matters.

\subsection{Honest disagreements}

We do not smooth over the cells where the cores differ. Archimedes is the
strongest of the three (12/15 against CONGA) and is the only core to reach DRILL
equivalence on the backbones, an advantage we attribute to its anticipatory,
threshold-based repulsion coping better with the backbones' sharp bottlenecks;
we report it but do not over-read a single-topology difference. Newton is the
weakest against CONGA (9/15), losing the smaller grids where its reactive scar
needs a drop before it can act. These differences are real and consistent with
each mechanism's character; none disturbs the central pattern of convergence.
